% ARCHIVED AUDITED DRAFT: incorporated into part_B.tex on 2026-11-19;
% independent root targeted proof audit PASS, 14:25 UTC.
\begin{saresult}{B.042}{Proved}{Every infinite Krull domain with arbitrary torsion-free coefficients}
Let $R$ be an infinite Krull domain with fraction field $K$, and
let $R\subset S$ be an $R$-torsion-free unital extension.
Then $\Omega(R,S)$ is a projective $R$-module, every
multivariable tensor map is an isomorphism, and the unique
compatible plethory exists. No Noetherian hypothesis is needed.
Here the Krull hypotheses used are precisely
$R=\bigcap_{\operatorname{ht}\mathfrak p=1}R_{\mathfrak p}$,
each $R_{\mathfrak p}$ a discrete valuation ring, and every
nonzero element of $R$ lying in only finitely many such primes.
\end{saresult}
\begin{proof}
The field case is B.014. Put $L=S_{\mathrm{rat}}\subset K$.
We verify the two places where B.020 used Noetherianity and
then complete the fractional-ideal decomposition directly.

Every height-one maximal ideal $\mathfrak p$ is invertible.
Choose $b\in R$ with $v_{\mathfrak p}(b)=1$. Besides
$\mathfrak p$, only finitely many height-one primes
$\mathfrak q_i$ contain $b$. Choose
$a_i\in\mathfrak q_i^{v_{\mathfrak q_i}(b)}\setminus\mathfrak p$
and put $a=\prod_i a_i$. Such choices exist because distinct
height-one primes are incomparable; powers of an element of
$\mathfrak q_i\setminus\mathfrak p$ suffice.
Then $x=a/b$ has valuation $-1$ at $\mathfrak p$ and
nonnegative valuation at every other height-one prime.
Consequently $x\mathfrak p\subset R$, by the Krull intersection
property. The ideal $\mathfrak p\mathfrak p^{-1}$ contains
$\mathfrak p$ and $xb=a\notin\mathfrak p$, so equals $R$.
This proves invertibility rather than assuming finite generation.

For each fixed prime power $q$, choose $r_q\in R$ with
$r_q^q-r_q\ne0$; this is possible since $R$ is infinite.
Every height-one maximal ideal with residue field of size $q$
contains this nonzero element, so there are only finitely many
such ideals. In particular the following product is finite:
\[
 I_n=\prod_{\substack{\operatorname{ht}\mathfrak p=1,\,
                       \mathfrak p\text{ maximal}\\
                       |R/\mathfrak p|=q<\infty}}
           \mathfrak p^{-w_q(n)},\qquad
 w_q(n)=\sum_{j\ge1}\lfloor n/q^j\rfloor.
\]
All its factors, and hence $I_n$, are invertible fractional ideals.
An invertible fractional ideal is the intersection of its
height-one localizations: if $x$ belongs to those localizations,
then $xI_n^{-1}$ is integral at every height-one prime and hence
lies in $R$, which implies $x\in I_n$.

The digit Newton argument at a finite-residue height-one maximal
ideal shows that the leading coefficient of an integer-valued
polynomial of degree at most $n$ has valuation at least
$-w_q(n)$. The necessary nodes can be chosen in $R$ by successive
residue digits and a valuation-one element in $R$.
At every other height-one prime, the image of $R$ in its residue
field is infinite, so Vandermonde interpolation gives
nonnegative valuations of all coefficients. The leading
coefficient ideal is therefore contained in $I_n$.

For each prime occurring in the finite product, choose its monic
local digit polynomial $Q_{n,\mathfrak p}\in R[X]$.
The Chinese remainder theorem for the pairwise comaximal
powers $\mathfrak p^{w_q(n)}$ gives a monic polynomial
$P_n\in R[X]$ of degree $n$, congruent coefficientwise to
$Q_{n,\mathfrak p}$ modulo that power at each such prime.
For $c\in I_n$ and $r\in R$, the product $cP_n(r)$ is integral
at every height-one prime, and so belongs to $R$.
Subtracting its leading term successively proves
\[
 \operatorname{Int}(R)=\bigoplus_{n\ge0}I_nP_n,\qquad
 \Omega(R,L)=\bigoplus_{n\ge0}(I_n\cap L)P_n.
\]
The second equality uses that the $P_n$ are monic with
coefficients in $R$, so descending elimination preserves $L[X]$.

Put $J_n=I_n\cap L$ and let $E_n$ be the finite set of primes
occurring in $I_n$. Localization commutes with this finite
intersection of submodules of $K$. At a maximal ideal outside
$E_n$, $J_n$ localizes to $R_{\mathfrak m}$. At
$\mathfrak p\in E_n$, the overring $L_{\mathfrak p}$ of the
rank-one discrete valuation ring $R_{\mathfrak p}$ is either
$R_{\mathfrak p}$ or $K$. Thus localization at every maximal
ideal proves the exact formula
\[
 J_n=\prod_{\substack{\mathfrak p\in E_n\\L_{\mathfrak p}=K}}
                 \mathfrak p^{-w_{|R/\mathfrak p|}(n)}.
\]
No finite-generation assertion about an arbitrary intersection
has been used: equality is proved by localization against the
displayed finite invertible product. Each $J_n$ is therefore
projective, and their direct sum is projective. B.019 proves
all tensor and plethory assertions.
\end{proof}

\begin{saresult}{B.043}{Proved}{Infinite-residue conductor cusps and all their coefficients}
Let $k=\mathbb F_q$ be finite, let $l\supset k$ be any infinite
field, and put
\[
 V=l[[t]],\qquad R=k+t^aV\quad(a\ge1),\qquad K=l((t)).
\]
For every $R$-torsion-free unital extension $R\subset S$, put
$L=S_{\mathrm{rat}}$. Then the polynomial coaddition and
comultiplication tensors are uniquely determined if and only
if $L=R$. In particular fraction-field coefficients fail.
For $A=\operatorname{Int}(R)$, $s\in l\setminus k$,
$Q(X)=X^q-X$ and $F(X)=Q(X)^2/t^a$, an explicit kernel is
\[
 (sF)\otimes F-F\otimes(sF)\ne0.
\]
The ring $R$ is a non-Noetherian local domain with maximal ideal
$t^aV$, residue field $k$, and fraction field $K$.
\end{saresult}
\begin{proof}
The ring assertions except non-Noetherianity follow from its
pullback description and from $t^aV\subset R$: an element with
nonzero constant in $k$ has inverse in $k+t^aV$, and division
by a power of $t^a$ gives every element of $K$.
Its maximal ideal $\mathfrak m=t^aV$ has
$\mathfrak m/\mathfrak m^2$ of infinite dimension over $k$,
since it contains $t^al$ and $[l:k]=\infty$.
Thus $\mathfrak m$ is not finitely generated.

We first treat coefficient field $K$ and construct a functional
which proves that the proposed tensor is nonzero.
For $G\in A$, the polynomial $G(t^aY)$ takes $V$ into
$R\subset V$. Since $V$ has infinite residue field,
$\operatorname{Int}(V)=V[Y]$ by Vandermonde interpolation.
Its reduction modulo $t$ is a polynomial over $l$ taking all
of $l$ into $k$. Such a polynomial is constant: the product of
its differences from all elements of $k$ vanishes at every
point of the infinite field $l$, hence is the zero polynomial;
one factor must be zero. Its constant is $G(0)\bmod t$.
It follows that
\[
 G(t^ay)-G(0)\in R\cap tV=t^aV\quad(y\in V).
\]
Apply the same Vandermonde argument to its quotient by $t^a$.
This defines an $R$-linear map, with $R$ acting through $k$,
\[
 h:A\longrightarrow l[Y],\qquad
 h(G)=\left(\frac{G(t^aY)-G(0)}{t^a}\right)\bmod t.
\]
In particular $h$ kills $\mathfrak m A$.

For $r\in R$ the value $Q(r)$ lies in $t^aV$, so both $F$ and
$sF$ take $R$ into $t^aV\subset R$. Direct expansion gives
\[
 h(F)=Y^2,\qquad h(sF)=sY^2.
\]
The image of the proposed tensor under $h\otimes h$ is
\[
 sY^2\otimes Z^2-Y^2\otimes sZ^2
       \in l[Y]\otimes_k l[Z],
\]
which is nonzero. For example choose a $k$-linear functional
$\ell:l\to k$ with $\ell(1)=0$, $\ell(s)=1$, and apply it
to the first $Y^2$ coefficient and take the second $Z^2$
coefficient. The result is $1$. The polynomial image of the
original tensor is zero by cancellation of its two products.

For completeness consider all coefficient overrings. Any
$L\not\subset V$ is $K$: if $v(b)<0$, choose $m$ with $v(b^{-m})\ge a$; then
$c=b^{-m}\in t^aV\subset R$ and $c^{-1}\in L$. Every element
of $K$ is sent into $R$ by a high power of $c$.
Suppose now $R\subsetneq L\subset V$. The conductor contains
$\mathfrak m$, and $Q(R)\subset\mathfrak m$.
Monic division gives a split summand $LQ$ in $\Omega(R,L)$.
If the second tensor map were injective, multiplication
$L\otimes_RL\to L$ would be an isomorphism. Modulo
$\mathfrak m$, it would give injective multiplication on the
nonzero $k$-algebra $C=L/\mathfrak mL=L/\mathfrak m$.
Choose a $k$-linear functional taking $1$ to $1$.
The relation $c\otimes1=1\otimes c$, forced by injectivity,
then implies that every $c\in C$ belongs to $k$.
Hence $C=k$ and $L=k+\mathfrak m=R$, a contradiction.
This argument does not use a finite-module Nakayama lemma:
$L$ can be infinite over $R$.
The only positive case is therefore $L=R$, which gives the
ordinary polynomial plethory. B.014 transfers the result to
arbitrary torsion-free $S$.
\end{proof}
